\begin{abstract}
	
	\textbf{Contexte.} Le réchauffement climatique modifie la structure des communautés phytoplanctoniques de l'océan Austral, avec une possible diminution de la contribution des diatomées au profit du nano- et picophytoplancton. Ce changement pourrait allonger le réseau trophique et affecter la biomasse des Niveaux Trophiques Intermédiaires (NTI, zooplancton et micronecton), avec des répercussions en cascade sur l'ensemble des écosystèmes.
	
	\vspace{0.3cm}
	
	\textbf{Objectif et méthode.} La relation entre la biomasse des NTI, estimée par acoustique active (Nautical Area Scattering Coefficient, NASC, à 38 et 120 kHz) a été modélisée avec des Forêts aléatoires par un ensemble de variables environnementales (température, salinité, hydrodynamisme via le Finite-Time Lyapunov Exponent) et pigmentaires (neuf pigments photosynthétiques). L'étude porte sur le secteur indien de l'océan Austral, en saison estivale (janvier-mars), à partir des campagnes OBSAUSTRAL 2018, 2021, 2022 et 2023. Une attention particulière a été portée à la fiabilité de l'évaluation et l'interprétation du modèle : une comparaison entre validation croisée Monte-Carlo et validation par blocage spatio-temporel a été effectuée. 
	%Des communautés phytoplanctoniques de la région ont ensuite été caractérisées par clustering k-means sur les concentrations pigmentaires. 
	
	\vspace{0.3cm}
	
	\textbf{Résultats.} Des cartes de prédictions de NASC ont été générées sur l'ensemble de la région pour chaque modèle de validation croisée. Malgré des écarts de performance importants entre les schémas de validation, les cartes obtenues sont similaires. Elles révèlent une structuration latitudinale inversée entre 38 et 120 kHz. Les domaines océanographiques fonctionnels (température/salinité) et la biomasse phytoplanctonique totale ressortent comme les plus importantes variables pour la prédiction du NASC et le microphytoplancton (diatomées, dinoflagellés) apparaît plus déterminant que le nano- et le picophytoplancton. La comparaison entre validation croisée Monte-Carlo et validation par blocage spatio-temporel révèle une fuite d'information substantielle dans le premier schéma, confirmée par une forte sensibilité au bruit.  Chaque schéma de validation témoigne d'un plafond structurel des courbes d'apprentissage dont l'origine, limite informationnelle des covariables, densité d'échantillonnage, ou non-stationnarité spatiale non prise en compte, n'a pu être totalement départagée. 
	%Une seconde analyse par clustering des concentrations pigmentaires met en évidence six communautés phytoplanctoniques structurées selon un gradient latitudinal sans toutefois révéler de différence significative de NASC entre ces communautés. 
	
	
	\vspace{0.3cm}
	
	\textbf{Conclusion et perspectives.} Ces résultats suggèrent que la structuration environnementale (physico-chimique et hydrodynamique) prédomine sur la seule composition pigmentaire dans l'explication de la distribution du NASC par le Random Forest, sur la période et l'échelle considérées. Ils soulignent également la nécessité d'une évaluation méthodologiquement rigoureuse dans les modèles écologiques spatio-temporels. Des travaux futurs pourraient intégrer une méthode de sélection de variable plus étoffée, la complétion de données manquantes, un décalage temporel entre efflorescence phytoplanctonique et réponse des NTI, l'intégration explicite de la non-stationnarité, l'exploitation de données de bio-logging pour élargir la couverture spatio-temporelle, et recourir à des modèles de machine learning causal pour identifier les relations biologiques entre les variables étudiées.
	
	\vspace{0.5cm}
	\noindent\textbf{Mots-clés :} océan Austral, niveaux trophiques intermédiaires, acoustique active, phytoplancton, arbres de régression, validation croisée spatio-temporelle
\end{abstract}